% ============================================================
\chapter{Functions of Random Variables}
\label{ch:functions_rv}
% ============================================================

In science and engineering, one rarely measures the quantity of direct interest. A physicist measures an angle and deduces a velocity; a financial analyst observes a log-return and deduces a price; an engineer measures a voltage and deduces a power. Each time, a transformation $g$ is applied to a random variable $X$ to obtain a new random variable $Y = g(X)$. But what is the distribution of $Y$? The question, innocent in appearance, leads to rich and varied techniques: the change-of-variable formula, the CDF method, Jacobians for multidimensional transformations. Mastering these tools means graduating from passive observer of randomness to architect capable of manipulating distributions at will.

\begin{intuition}
When a transformation is applied to a random variable, the result is a new random
variable whose distribution may be very different. This chapter develops the methods
for determining the distribution of $Y = g(X)$ from that of $X$.
\end{intuition}

% ------------------------------------------------------------
\section{Discrete Case}
% ------------------------------------------------------------

\begin{proposition}[Distribution of $g(X)$ --- discrete case]
If $X$ is discrete taking values $x_1, x_2, \ldots$ and $g : \R \to \R$, then
$Y = g(X)$ is discrete with
\[
    \PP(Y = y) = \sum_{k:\, g(x_k)=y} \PP(X = x_k).
\]
\end{proposition}

\begin{example}
If $X \sim \mathrm{Bin}(n, 1/2)$ and $Y = 2X - n$, then $Y$ takes values
$-n, -n+2, \ldots, n-2, n$ and
\[
    \PP(Y = 2k-n) = \binom{n}{k}\left(\frac{1}{2}\right)^n, \quad k = 0, 1, \ldots, n.
\]
\end{example}

% ------------------------------------------------------------
\section{The CDF Method}
% ------------------------------------------------------------

\begin{theorem}[CDF method]
\label{thm:cdf_method}
To find the distribution of $Y = g(X)$:
\begin{enumerate}
    \item Compute $F_Y(y) = \PP(Y \leq y) = \PP(g(X) \leq y)$.
    \item Differentiate: $f_Y(y) = F_Y'(y)$.
\end{enumerate}
\end{theorem}

\begin{example}[$Y = X^2$ with $X \sim \mathcal{N}(0,1)$]
For $y > 0$:
\begin{align*}
    F_Y(y) &= \PP(X^2 \leq y) = \PP(-\sqrt{y} \leq X \leq \sqrt{y})
    = \Phi(\sqrt{y}) - \Phi(-\sqrt{y}) = 2\Phi(\sqrt{y}) - 1.
\end{align*}
Differentiating:
\[
    f_Y(y) = 2 \cdot \frac{1}{\sqrt{2\pi}}e^{-y/2}\cdot\frac{1}{2\sqrt{y}}
    = \frac{1}{\sqrt{2\pi}}\,y^{-1/2}\,e^{-y/2}, \quad y > 0.
\]
This is the density of $\Gamma(1/2, 1/2) = \chi^2(1)$.
\end{example}

\begin{example}[$Y = e^X$ with $X \sim \mathcal{N}(\mu, \sigma^2)$]
For $y > 0$:
\[
    F_Y(y) = \PP(e^X \leq y) = \PP(X \leq \ln y)
    = \Phi\!\left(\frac{\ln y - \mu}{\sigma}\right).
\]
Differentiating:
\[
    f_Y(y) = \frac{1}{y\sigma\sqrt{2\pi}}
    \exp\!\left(-\frac{(\ln y - \mu)^2}{2\sigma^2}\right), \quad y > 0.
\]
This is the log-normal density.
\end{example}

% ------------------------------------------------------------
\section{The Change-of-Variable Formula}
% ------------------------------------------------------------

\begin{theorem}[Change of variable --- monotone case]
\label{thm:change_var}
Let $X$ be a continuous variable with density $f_X$ and $g$ a
\textbf{strictly monotone} $C^1$ function on the support of $X$, with $C^1$
inverse $g^{-1}$. Then $Y = g(X)$ has density
\[
    f_Y(y) = f_X\bigl(g^{-1}(y)\bigr)\cdot\abs{(g^{-1})'(y)}.
\]
\end{theorem}

\begin{proof}
Suppose $g$ is strictly increasing (the decreasing case is analogous).
\[
    F_Y(y) = \PP(g(X) \leq y) = \PP(X \leq g^{-1}(y)) = F_X(g^{-1}(y)).
\]
Differentiating by the chain rule:
$f_Y(y) = f_X(g^{-1}(y))\cdot (g^{-1})'(y)$.
Since $g$ is increasing, $(g^{-1})' > 0$, so the absolute value is unnecessary.
In the decreasing case a minus sign appears, absorbed by $\abs{\cdot}$.
\end{proof}

\begin{example}[Affine transformation]
If $Y = aX + b$ with $a \neq 0$:
$g^{-1}(y) = (y-b)/a$, $(g^{-1})'(y) = 1/a$. Hence:
\[
    f_Y(y) = \frac{1}{\abs{a}}\,f_X\!\left(\frac{y-b}{a}\right).
\]
In particular, if $X \sim \mathcal{N}(\mu, \sigma^2)$, then
$Z = (X-\mu)/\sigma \sim \mathcal{N}(0,1)$.
\end{example}

% ------------------------------------------------------------
\section{Non-Monotone Transformations}
% ------------------------------------------------------------

\begin{theorem}[General change of variable]
If $g$ is $C^1$ and the support of $X$ can be decomposed into finitely many
intervals $I_1, \ldots, I_m$ on each of which $g$ is strictly monotone, with
inverses $h_1, \ldots, h_m$, then
\[
    f_Y(y) = \sum_{j=1}^{m} f_X(h_j(y))\cdot\abs{h_j'(y)}\,\mathbf{1}_{g(I_j)}(y).
\]
\end{theorem}

\begin{example}[$Y = X^2$ with symmetric $f_X$]
$g(x) = x^2$ is decreasing on $(-\infty, 0)$ and increasing on $(0, +\infty)$.
The two inverses are $h_1(y) = -\sqrt{y}$ and $h_2(y) = \sqrt{y}$, with
$\abs{h_j'(y)} = 1/(2\sqrt{y})$. Hence:
\[
    f_Y(y) = \frac{f_X(\sqrt{y}) + f_X(-\sqrt{y})}{2\sqrt{y}}, \quad y > 0.
\]
\end{example}

% ------------------------------------------------------------
\section{Moment Generating Functions}
% ------------------------------------------------------------

\begin{definition}[Moment generating function (MGF)]
The \textbf{moment generating function} of $X$ is
\[
    M_X(t) = \E[e^{tX}],
\]
defined for those $t$ in a neighbourhood of $0$ where this expectation exists.
\end{definition}

\begin{proposition}[Properties of the MGF]
\begin{enumerate}[label=(\roman*)]
    \item $M_X(0) = 1$.
    \item $M_X^{(n)}(0) = \E[X^n]$ (moments are obtained by differentiation).
    \item If $M_X(t) = M_Y(t)$ for all $t$ in a neighbourhood of $0$, then $X$
          and $Y$ have the same distribution.
    \item If $X \indep Y$, then $M_{X+Y}(t) = M_X(t)\cdot M_Y(t)$.
\end{enumerate}
\end{proposition}

\begin{example}[MGF of the normal distribution]
If $X \sim \mathcal{N}(\mu, \sigma^2)$:
\[
    M_X(t) = \exp\!\left(\mu t + \frac{\sigma^2 t^2}{2}\right).
\]
We recover $M_X'(0) = \mu = \E[X]$ and $M_X''(0) = \sigma^2 + \mu^2 = \E[X^2]$.
\end{example}

% ------------------------------------------------------------
\section{Probability Integral Transform}
% ------------------------------------------------------------

\begin{theorem}[Probability integral transform]
If $X$ is a continuous variable with strictly increasing CDF $F_X$, then
$U = F_X(X) \sim \mathcal{U}([0,1])$.
\end{theorem}

\begin{proof}
For $u \in (0,1)$:
$\PP(U \leq u) = \PP(F_X(X) \leq u) = \PP(X \leq F_X^{-1}(u))
= F_X(F_X^{-1}(u)) = u$.
\end{proof}

\begin{remark}
This result is the theoretical basis for the inverse transform simulation method
(Chapter~4) and for goodness-of-fit tests ($p$-values).
\end{remark}

% ------------------------------------------------------------
\section{Exercises}
% ------------------------------------------------------------

\begin{exercise}
Let $X \sim \mathcal{U}([0,1])$ and $Y = -\ln X$. Find the distribution of $Y$.
\end{exercise}

\begin{exercise}
Let $X \sim \mathrm{Exp}(1)$. Find the density of $Y = \sqrt{X}$.
\end{exercise}

\begin{exercise}
Let $X \sim \mathcal{N}(0,1)$. Find the density of $Y = \abs{X}$
(half-normal distribution).
\end{exercise}

\begin{exercise}
Let $X$ have density $f_X(x) = 2x\,\mathbf{1}_{[0,1]}(x)$ and $Y = 1 - X^2$.
Find the density of $Y$.
\end{exercise}

\begin{exercise}
Compute the MGF of $\mathrm{Exp}(\lambda)$ and deduce $\E[X]$, $\E[X^2]$, $\Var(X)$.
\end{exercise}

\begin{exercise}
Let $X \sim \mathrm{Cauchy}(0,1)$ with density $f(x) = 1/(\pi(1+x^2))$.
Show that the MGF of $X$ does not exist (i.e.\ $\E[e^{tX}] = \infty$ for
all $t \neq 0$).
\end{exercise}

\begin{exercise}[Box--Muller]
Let $U_1, U_2$ be i.i.d.\ $\mathcal{U}([0,1])$. Define
$Z_1 = \sqrt{-2\ln U_1}\cos(2\pi U_2)$ and
$Z_2 = \sqrt{-2\ln U_1}\sin(2\pi U_2)$.
Show that $Z_1, Z_2$ are i.i.d.\ $\mathcal{N}(0,1)$.
\end{exercise}

\begin{exercise}
Using the MGF, show that if $X_1 \sim \mathcal{N}(\mu_1, \sigma_1^2)$ and
$X_2 \sim \mathcal{N}(\mu_2, \sigma_2^2)$ are independent, then
$X_1 + X_2 \sim \mathcal{N}(\mu_1+\mu_2, \sigma_1^2+\sigma_2^2)$.
\end{exercise}
